
\section{Pair production by a photon in the field of a nucleus}

\SourcePage{454}{459}

The production of an electron-positron pair in the collision of a photon with a nucleus ($Z + \gamma \to Z + e^- + e^+$) and bremsstrahlung in the collision of an electron with a nucleus ($Z + e^- \to Z + e^+\gamma$) are two crossed channels of one and the same reaction. In \S\,91 the rules were already formulated by which the formulas are transformed in passing from the second of these cases to the first. In the present case, applying these rules to formula (93.8), we obtain the following expression for the differential cross section of pair production by an unpolarised photon, averaged over the polarisations of the pair components\footnote{On polarisation effects in pair production by photons see the same literature as indicated in \S\,93 for bremsstrahlung.}:

\SourcePage{455}{460}

\begin{equation}
\begin{split}
d\sigma &= \frac{Z^2\alpha r_\mathrm{e}^2\,m^2 p_+ p_-}{2\pi\,\omega^3 q^4}\sin\theta_+\,d\theta_+\sin\theta_-\,d\theta_-\,d\varphi \times {} \\
&\quad \times \biggl\{\frac{p_+^2}{\varkappa_+^2}(4\varepsilon_-^2 - q^2)\sin^2\theta_+ + \frac{p_-^2}{\varkappa_-^2}(4\varepsilon_+^2 - q^2)\sin^2\theta_- - {} \\
&\quad {} - \frac{2\omega^2}{\varkappa_+\varkappa_-}(p_+^2\sin^2\theta_+ + p_-^2\sin^2\theta_-) - {} \\
&\quad {} - \frac{2p_+p_-}{\varkappa_+\varkappa_-}(2\varepsilon_+^2 + 2\varepsilon_-^2 - q^2)\sin\theta_+\sin\theta_-\cos\varphi\biggr\},
\end{split}
\label{eq:94.1}
\end{equation}
\[
\varkappa_\pm = \varepsilon_\pm - p_\pm\cos\theta_\pm, \quad q^2 = (\mathbf{p}_+ + \mathbf{p}_- - \mathbf{k})^2, \quad \varepsilon_+ + \varepsilon_- = \omega
\]
(\emph{H.~A.~Bethe, W.~Heitler}, 1934).

By the same transformation we obtain from (93.9) the distribution of the pair components in energy:
\begin{equation}
\begin{split}
d\sigma_{\varepsilon_+} &= Z^2\alpha r_\mathrm{e}^2\,\frac{p_+p_-}{\omega^3}\,d\varepsilon_+\biggl\{-\frac{4}{3} - 2\varepsilon_+\varepsilon_-\,\frac{p_+^2 + p_-^2}{p_+^2 p_-^2} + {} \\
&\quad {} + m^2\!\left(l_-\frac{\varepsilon_+}{p_-^2} + l_+\frac{\varepsilon_-}{p_+^2} - \frac{l_+l_-}{p_+p_-}\right) + {} \\
&\quad {} + L\biggl[-\frac{8\varepsilon_+\varepsilon_-}{3p_+p_-} + \frac{\omega^3}{p_+^3 p_-^3}\bigl(\varepsilon_+^2\varepsilon_-^2 + p_+^2 p_-^2 - m^2\varepsilon_+\varepsilon_-\bigr) - {} \\
&\quad {} - \frac{m^2\omega}{2p_+p_-}\!\left(l_+\frac{\varepsilon_+\varepsilon_- - p_+^2}{p_+^3} + l_-\frac{\varepsilon_+\varepsilon_- - p_-^2}{p_-^3}\right)\biggr]\biggr\},
\end{split}
\label{eq:94.2}
\end{equation}
\[
L = \ln\frac{\varepsilon_+\varepsilon_- + p_+p_- + m^2}{\varepsilon_+\varepsilon_- - p_+p_- + m^2}, \quad l_\pm = \ln\frac{\varepsilon_\pm + p_\pm}{\varepsilon_\pm - p_\pm}.
\]

\SourcePage{456}{461}

In the ultrarelativistic case ($\varepsilon_\pm \gg m$) the electron and positron emerge at angles $\theta_\pm \sim m/\varepsilon_\pm$ to the direction of the incident photon. The angular distribution is given by a formula analogous to (93.13):
\begin{equation}
\begin{split}
d\sigma &= \frac{8}{\pi}\,Z^2\alpha r_\mathrm{e}^2\,\frac{m^2\varepsilon_+\varepsilon_-}{\omega^3 q^4}\,d\varepsilon_+ \times {} \\
&\quad \times \biggl\{-\frac{\delta_+^2}{(1+\delta_+^2)^2} - \frac{\delta_-^2}{(1+\delta_-^2)^2} + \frac{\omega^2}{2\varepsilon_+\varepsilon_-}\,\frac{\delta_+^2 + \delta_-^2}{(1+\delta_+^2)(1+\delta_-^2)} + {} \\
&\quad {} + \left(\frac{\varepsilon_+}{\varepsilon_-} + \frac{\varepsilon_-}{\varepsilon_+}\right)\frac{\delta_+\delta_-\cos\varphi}{(1+\delta_+^2)(1+\delta_-^2)}\biggr\}\delta_+\delta_-\,d\delta_+\,d\delta_-\,d\varphi,
\end{split}
\label{eq:94.3}
\end{equation}
where
\begin{equation}
\frac{q^2}{m^2} = \delta_+^2 + \delta_-^2 - 2\delta_+\delta_-\cos\varphi + m^2\!\left(\frac{1+\delta_+^2}{2\varepsilon_+} + \frac{1+\delta_-^2}{2\varepsilon_-}\right)^{\!2}.
\label{eq:94.4}
\end{equation}

The energy distribution in this case:
\begin{equation}
d\sigma = 4Z^2\alpha r_\mathrm{e}^2\,\frac{d\varepsilon_+}{\omega^3}\left(\varepsilon_+^2 + \varepsilon_-^2 + \frac{2}{3}\,\varepsilon_+\varepsilon_-\right)\!\left(\ln\frac{2\varepsilon_+\varepsilon_-}{m\omega} - \frac{1}{2}\right) \quad \text{(u.r.)}.
\label{eq:94.5}
\end{equation}

Integration of this expression over $\varepsilon_+$ (in the limits from $m$ to $\omega$) gives the total cross section of pair production by a photon of given energy\footnote{Because of the convergence of the integral at both limits, the inapplicability of formula (94.5) when $\varepsilon_\pm - m$ is small is inessential.}:
\begin{equation}
\sigma = \frac{28}{9}\,Z^2\alpha r_\mathrm{e}^2\!\left(\ln\frac{2\omega}{m} - \frac{109}{42}\right), \quad \omega \gg m.
\label{eq:94.6}
\end{equation}

As for bremsstrahlung, the logarithmic term in the cross section in the ultrarelativistic case originates from the region of values $q \sim m^2/\varepsilon$. To this there now correspond angles for which
\[
|\delta_+ - \delta_-| \lesssim \frac{m}{\varepsilon}, \quad |\pi - \varphi| \lesssim \frac{m}{\varepsilon}
\]
(instead of $\varphi \lesssim m/\varepsilon$ in (93.15)). Thus, in the logarithmic approximation the directions of the electron and positron make small angles inversely proportional to the particle energies with the direction of the photon, and lie almost in one plane with the latter but on opposite sides of it.

Near the reaction threshold ($\omega \to 2m$) the Born approximation is inapplicable. The derivation of a quantitative formula in this case would require exact account of the Coulomb interaction of the three charged particles present in the final state (the nucleus and the pair). The symmetry with respect to the electron (attracted

\SourcePage{457}{462}

to the nucleus) and the positron (repelled from the nucleus) in this case, of course, disappears.

If
\begin{equation}
Z\alpha\sqrt{\frac{\omega - 2m}{\omega}} \ll 1,
\label{eq:94.7}
\end{equation}
then the Born approximation is still applicable. At non-relativistic energies of the pair $\omega \approx 2m \gg p_\pm$, and therefore $q \approx \omega$. In (94.1) one can put everywhere $\varepsilon_\pm = \varkappa_\pm = m$, $\omega = 2m$, after which this formula reduces to the expression
\begin{equation}
d\sigma = \frac{Z^2\alpha r_\mathrm{e}^2\,p_+p_-}{64\pi^2\,m^6}(p_+^2\sin^2\theta_+ + p_-^2\sin^2\theta_-)\,do_+\,do_-\,d\varepsilon_+.
\label{eq:94.8}
\end{equation}

After integration over angles
\begin{equation}
\begin{split}
d\sigma &= \frac{1}{6}\,Z^2\alpha r_\mathrm{e}^2\,\frac{p_+p_-(p_+^2 + p_-^2)}{m^5}\,d\varepsilon_+ = {} \\
&= \frac{2Z^2\alpha r_\mathrm{e}^2}{3m^3}\,(\omega - 2m)\sqrt{(\varepsilon_+ - m)(\varepsilon_- - m)}\,d\varepsilon_+.
\end{split}
\label{eq:94.9}
\end{equation}

Finally, integrating over $\varepsilon_+$ (in the limits from $m$ to $\omega - m$), we obtain the total cross section
\begin{equation}
\sigma = \frac{\pi}{12}\,Z^2\alpha r_\mathrm{e}^2\!\left(\frac{\omega - 2m}{m}\right)^{\!3}.
\label{eq:94.10}
\end{equation}

If the relative velocity $v_0$ of the components of the produced pair is small, it is necessary to take into account their Coulomb interaction with each other (\emph{A.~D.~Sakharov}, 1948). It becomes substantial when $v_0$ is of the order of (or less than) the velocity of a particle in the bound state of the electron and positron (positronium):
\begin{equation}
v_0 \lesssim \alpha.
\label{eq:94.11}
\end{equation}

Let us consider the process in the centre-of-mass frame of the pair. In the diagrams depicting the process in this frame, the essential virtual momenta are $\sim m$. In other words, the essential distances between the electron and positron are $\sim 1/m$. Meanwhile the wave function of their relative motion $\psi(r)$ changes substantially only at distances $r \sim 1/(mv_0) \sim 1/(m\alpha)$, i.e.\ large compared with $1/m$. Therefore the account of the interaction of the particles reduces to the appearance in the transition matrix element of a factor $\psi^*(0)$. Correspondingly, the differential cross section is multiplied by $|\psi(0)|^2$, i.e.\ by
\begin{equation}
\frac{2\pi\alpha/v_0}{1 - e^{-2\pi\alpha/v_0}}
\label{eq:94.12}
\end{equation}

\SourcePage{458}{463}

(see III, (136.11)). The relative velocity of two particles is the velocity of one of them in the rest frame of the other. Comparing the values of the invariant $p_{-\mu}p_+^\mu$ in this frame and in the laboratory frame (the rest frame of the nucleus), we obtain
\[
\frac{m^2}{\sqrt{1 - v_0^2}} = \varepsilon_+\varepsilon_- - \mathbf{p}_+\mathbf{p}_-,
\]
from which $v_0$ can be found. If $\mathbf{p}_+$ and $\mathbf{p}_-$ are close to each other in magnitude and direction, then for $v_0$ we obtain the approximate formula
\begin{equation}
v_0^2 = \frac{p^2}{m^2}\,\vartheta^2 + \frac{(p_+ - p_-)^2}{\varepsilon^2},
\label{eq:94.13}
\end{equation}
applicable when $v_0 \ll 1$; $p = (p_+ + p_-)/2$, $\varepsilon = (\varepsilon_+ + \varepsilon_-)/2$, $\vartheta$ is the angle between $\mathbf{p}_+$ and $\mathbf{p}_-$.

The correction to the cross section, determined by formulas (94.12), (94.13), leads to the appearance of an anomaly in the correlation between the momenta of the produced electron and positron: a narrow maximum at $\mathbf{p}_+ \approx \mathbf{p}_-$.
